\section{问题一的模型建立与求解}
% 写作提示：围绕问题目标依次交代数据与输入输出、模型准备、变量和参数、核心关系式或目标与约束、
% 方法选择依据、算法步骤与参数依据、计算环境、关键结果及其解释。下方模型和数值仅为示例，须按实际计算替换。
\subsection{数据预处理}
采用Z-score标准化处理：$x' = \frac{x - \mu}{\sigma}$。

\subsection{模型建立}
采用多元线性回归作为基础模型：
\begin{equation}
y = \theta_0 + \theta_1 x_1 + \theta_2 x_2 + \cdots + \theta_n x_n + \varepsilon
\end{equation}
同时引入随机森林模型进行对比分析，通过集成多棵决策树提高精度和鲁棒性。

\subsection{求解结果}
随机森林的测试集$R^2$达到0.896，优于线性回归的0.834。
